\section{Three ways to search the selected documents}
\subsection{A: cluster, then full scan}
We first use a scan because its behavior is easy to check. In the six-document example,
coordinate 1 defines two clusters. We open cluster 1 and then cluster 0. A reads the codes
of documents 2, 3, 4, then 0, 1, 5, computes their full scores, and keeps the best two.
The order of visits does not change the final ranking.

A \emph{top-$K$ heap} stores the best $K$ results found so far, with the worst retained result
available for replacement. For $K=2$, after documents 2 and 3 it holds scores 0.1 and 2.5.
Document 4, scoring 1.9, replaces document 2. None of the remaining documents improves the pair.

\begin{lstlisting}[caption={Full scan. Every selected document receives a full score.}]
scan(query, selected_clusters, K):
    table = prepare_score_table(query)
    best = empty_top_k(K)
    for cluster in selected_clusters:
        for document in cluster.documents:
            score = full_score(table, document.code)
            best.offer(score, document.id)
    return best.sorted_by_score_then_id()
\end{lstlisting}

The score table avoids multiplying every sign separately. For each group of four coordinates,
we precompute the contribution of all 16 possible sign patterns. A 256-bit document then needs
64 table terms. A lookup still involves reading bits, extracting a pattern, loading a value,
and adding it to the score. A table term is not a single CPU instruction.

More generally, for $g$ coordinates per score-table group, there are
$G=\ceil{d/g}$ groups and $G2^g$ stored table values, allowing padding in the final group.
Preparing the table is query work, shared by all three methods. Changing $g$ trades table size
against the number of terms per document. The measured implementation uses $g=4$.

Let $\mathcal P$ be the selected clusters and $n_c$ the size of cluster $c$. A scores
\[
F_A=\sum_{c\in\mathcal P}n_c
\]
documents and evaluates $G F_A$ table terms. A useful local time model is
\begin{equation}
T_A^{\rm local}\approx T_{\rm table}+F_Ac_{\rm scan}+T_{\rm select}(F_A,K).
\label{eq:scan-time}
\end{equation}
Here $c_{\rm scan}$ is the cost of reading and scoring one contiguous document at the chosen
width. Selection is written separately to show its role. If a measured per-row coefficient
already includes selection, we must not add it again. A worst-case heap bound is
$O(F_A\log K)$, but not every score changes the heap.

\subsection{The scan's storage and temporary memory}
Let $w$ be the number of bits in a packed word and $u=w/8$ its bytes. A code occupies
$u\ceil{d/w}$ bytes; let $b_{\rm id}$ be the document-ID width. The row payload is
\begin{equation}
M_{\rm rows}=N\left(u\ceil{d/w}+b_{\rm id}\right).
\label{eq:rows}
\end{equation}
For $d=256$, $w=64$, and eight-byte IDs, each row needs $32+8=40$ bytes. At one million rows,
that is 40 MB in decimal units. Cluster directories, vector spare capacity, and allocator
bookkeeping require additional memory.

Routing has its own storage. For example, $C$ float centroids of width $d_r$ need $4Cd_r$
bytes for their coordinates. Exact centroid comparison examines $Cd_r$ coordinates per query,
plus the work of choosing the best $P$ centroids. A sign-prefix router has a different cost.
Routing cannot be assigned one universal price independent of its implementation.

The query also needs temporary storage: its vector, the score table, selected cluster IDs,
and the current results. Denote this shared allocation by $Q_0$. Its main terms scale as
$O(d+G2^g+P+K)$. Include Python/native result arrays if they coexist. The serving memory for A is
stored rows, routing and container allocation, plus $Q_0$. A build can peak higher because the
input arrays and newly arranged index exist together. File size, logical payload, and process
RAM therefore answer different questions.

The scan has a useful hardware property: codes in a cluster are read consecutively. The CPU
can fetch future cache lines while it works on the current ones. SIMD instructions can process
several values together. These effects change the measured price of a row, even though the
number of scored documents is unchanged. They explain why reducing arithmetic alone is not
enough to demonstrate a faster search.

\subsection{B: arrange the same bits by coordinate}
To avoid some full scores, B needs a way to select documents by a query bit. It stores a
\emph{bitplane}: one coordinate across all documents in a cluster. A bitplane is the same sign
data turned into columns. B keeps the packed document rows as well, so it can score a selected
candidate with the common full scorer.

In cluster 1 of our example, the local document order is [2, 3, 4]. Their codes are
\texttt{10100}, \texttt{11001}, and \texttt{11000}. The plane for coordinate 2 is
\texttt{011}: documents 3 and 4 have that bit set. The plane for coordinate 3 is
\texttt{100}. The query prefers coordinate 2 to be 1 and coordinate 3 to be 0; both selections
keep documents 3 and 4.

A candidate mask records which local document positions remain eligible. With parent mask
$A$ and plane $B_i$, a positive-query split makes
\[
A_{\rm match}=A\mathbin{\&}B_i,
\qquad A_{\rm other}=A\mathbin{\&}\neg B_i.
\]
For a negative query coordinate, the matching and other children swap. Each result is
intersected with the parent, so padding positions never become real documents.

\fig{masks}{Membership shrinks, but a dense mask keeps its physical width. Each child still has 6,400 positions, or 100 words. This larger illustration counts storage; the six-document example explains the returned IDs.}{.94}

\subsection{Follow the promising branch, keep the alternative}
The query's magnitudes order coordinates as 1, 4, 2, 3, 5. Within cluster 1, coordinate 1 is
already fixed and all three documents agree with the query on coordinate 4. Splitting on
coordinate 2 then separates [3, 4] from [2]. The first child has no new mismatch penalty;
the second adds 0.5. We visit the more promising child first.

Visiting it first does not justify deleting the alternative. A document may match this bit
but disagree on several later bits. B stores unexplored children in a queue, ordered by a
lower bound on their mismatch penalty. From Equation~\ref{eq:score}, a node with known
penalty $p$ has score upper bound $Q-2p$: unexamined coordinates are allowed their best signs.
If this upper bound is strictly below the worst retained top-$K$ score, the node cannot help.
Equal-score possibilities still need the ID tie rule; the native check is conservative around
floating-point comparisons.

\begin{lstlisting}[caption={Bitplane search. The queue preserves alternatives to the preferred path.}]
branch(query, clusters, K, leaf_size, node_budget):
    table = prepare_score_table(query)
    order = coordinates_by_decreasing_abs(query)
    queue = one full candidate mask per selected cluster
    best = empty_top_k(K)
    while queue is not empty:
        node = remove_lowest_penalty(queue)
        if cannot_improve(node, query, best): continue
        if node_budget_is_reached(): return best
        if node.count <= leaf_size or no_bits_remain(node):
            score_each_set_document(table, node.mask, best)
        else:
            matching, other = split_on_next_bit(node, order)
            queue_nonempty_children(queue, matching, other)
    return best.sorted_by_score_then_id()
\end{lstlisting}

This listing shows the search policy; the native source also handles counters, ownership,
and numerical tolerance. A leaf is a node small enough to score directly. The leaf size $L$
therefore decides when to stop splitting. A node budget limits the work but may stop before
the bounds establish the answer. With no work limit and safe bounds, the method can finish
exactly within the selected clusters. Unprobed clusters are still outside that guarantee.

A probability parameter can occasionally choose another queued node instead of the best bound.
This changes the order of exploration, not the score. Under a fixed budget it may help or hurt
recall, so we test it with seeds rather than assume an improvement. The tested real-data
settings did not show a consistent gain from adding random exploration.

\subsection{Count the full-width work}
For cluster $c$, each dense mask spans $W_c=\ceil{n_c/w}$ words. Let $J_c$ count its splits
and $E_c$ the leaf masks whose set bits are enumerated. The full-width visits are
\begin{equation}
V_{\rm split}=\sum_{c\in\mathcal P}J_cW_c,
\qquad V_{\rm leaf}=\sum_{c\in\mathcal P}E_cW_c.
\label{eq:visits}
\end{equation}
A leaf loop still checks zero words to find out that they are empty. The number of set bits
then determines how many documents are scored. Omitting leaf-word visits was one source of
an overly optimistic early cost estimate.

Suppose a 6,400-document cluster has been narrowed to 80 candidates. Its mask still contains
100 words. Ten further dense splits visit 1,000 words, even if later masks contain only a few
set bits. Each visit reads the parent and plane, builds children, and writes results. Allocating
masks and maintaining the queue cost time as well.

If $F_B$ documents are eventually scored, a practical decomposition is
\begin{equation}
\begin{split}
T_B^{\rm local}\approx{}&T_{\rm table}+T_{\rm order}
 +V_{\rm split}c_{\rm split}+V_{\rm leaf}c_{\rm leaf}\\
 &+F_Bc_{\rm gather}+T_{\rm queue}+T_{\rm select}(F_B,K).
\end{split}
\label{eq:branch-time}
\end{equation}
The per-word split price should include the chosen mask allocation policy. Scattered document
reads may have a different per-row price from contiguous scan reads. If a measured leaf timer
already includes finding set bits, fetching rows, and scoring them, use that combined term
instead of adding these stages a second time.

B's additional stored plane payload is
\begin{equation}
M_{\rm planes}=du\sum_{c=1}^{C}\ceil{n_c/w}.
\label{eq:planes}
\end{equation}
For one large cluster this is approximately $Nd/8$ bytes. Many small clusters add padding.
During the query, if $A_c(t)$ masks for cluster $c$ coexist at time $t$, the mask payload peak is
$u\max_t\sum_c A_c(t)W_c$. Add queue-node allocation and shared scratch. Counting just the
currently visited mask misses the alternatives waiting in the queue.

\subsection{C: begin with the query's key}
B chooses a sequence of columns from the query. C instead builds a directory over a fixed set
of $h$ coordinates. Within each cluster, documents with the same short key are stored together.
A directory entry maps an occupied key to a start position and count. The existing document IDs
identify that posting range; the implementation does not add a second ID array for postings.

Use coordinates 2 and 3 as the key in the running example. The query prefers \texttt{10}.
Flipping coordinate 3 costs 0.4, flipping coordinate 2 costs 0.5, and flipping both costs 0.9.
Thus the weighted order is \texttt{10}, \texttt{11}, \texttt{00}, \texttt{01}.

\fig{keys}{One shared key sequence is looked up in both selected clusters. The ideal key already finds the top two, but that fact alone is not a safe stopping rule. Empty lookups still cost work.}{.98}

The ideal key returns document 0 from cluster 0 and documents 3 and 4 from cluster 1. Their
full scores are 1.1, 2.5, and 1.9, so the best pair is [3, 4]. The next key's known penalty is
0.4, giving a full-score upper bound of $2.5-2(0.4)=1.7$. That is below 1.9, so in this
example a safe bound can stop before the empty \texttt{11} lookups. Figure~\ref{fig:keys} shows the wider
sequence to make the directory and empty keys visible; an unlimited walk without that bound
would inspect them too.

Why not always stop after finding $K$ documents? A short key ignores the other coordinates.
A document matching the key may have a poor full score, while another key may contain a better
document. C therefore scores the full code and keeps an upper bound for unvisited keys.

\begin{lstlisting}[caption={Weighted key lookup. Generate the key sequence once; charge every cluster lookup.}]
keys(query, clusters, K, candidate_target, lookup_limit):
    table = prepare_score_table(query)
    ideal = signs_of_fixed_key_coordinates(query)
    queue = [empty_flip_set_with_zero_penalty]
    best = empty_top_k(K)
    while queue is not empty:
        if remaining_bound_cannot_improve(queue, query, best): break
        state = remove_lowest_penalty(queue)
        key = ideal XOR state.flip_mask
        for cluster in clusters:
            if lookup_limit_is_reached(): return best
            count_one_lookup_including_misses()
            for document in cluster.directory.get(key, []):
                best.offer(full_score(table, document.code), document.id)
        if candidate_target_is_reached(): return best
        add_next_unique_flip_sets(queue, state)
    return best.sorted_by_score_then_id()
\end{lstlisting}

The queue enumerates combinations in increasing total mismatch weight. It is not restricted
to increasing numbers of flipped bits: two low-weight flips can cost less than one high-weight
flip. The native implementation uses a subset-generation rule that avoids generating the same
combination repeatedly. For every unvisited key, $Q-2p_{\rm key}$ is a safe full-score upper
bound, because omitted coordinates may all match. That bound can be loose.

A positive candidate target or lookup limit can stop early and lose recall; zero disables a
limit in the native API. The target is checked after completing a key across selected clusters.
The lookup limit may stop during that cluster loop. Exact local completion requires exhausting
potentially improving keys or proving that none remain, while respecting score ties.

\subsection{What the key directory costs}
Let $H$ count actual directory lookups and $F_C$ the full scores. With $E$ generated keys and
no limit cutting a cluster loop, $H=PE$. An empty key contributes to $H$ but not $F_C$. A time
model for the implemented policy is
\begin{equation}
\begin{split}
T_C^{\rm local}\approx{}&T_{\rm table}+T_{\rm key\ order}+T_{\rm key\ queue}
 +Hc_{\rm lookup}\\
 &+F_Cc_{\rm posting}+T_{\rm select}(F_C,K).
\end{split}
\label{eq:key-time}
\end{equation}
Posting rows are contiguous within a key, but many short ranges still differ from one long
scan. Queue costs depend on the live generated states. We must not charge $P$ copies of the
queue: this implementation shares one stream across clusters.

For $U_c$ occupied keys in cluster $c$, the logical directory payload is
$\sum_c U_c(b_k+2b_o)$: a key, start, and count. The current key field uses four bytes and each
range field uses eight. Allocated memory is larger. The C++ directory uses separate chaining,
with one node per occupied entry and a bucket-pointer array. If $B_c$ is the number of hash
buckets, its allocation includes $\sum_c(U_cb_{\rm node}+B_cb_{\rm ptr})$, plus container and
allocator overhead. The 20-byte field sum is not an allocated-node measurement.

A wider key can shorten posting ranges while making useful keys harder to reach. A very short
key finds documents easily but may leave nearly all the scoring work. Finding a nonempty key
and finding the best documents are different events. Looking up a relaxed pattern such as
\texttt{1*0} also needs an explicitly different operation: a hash table over exact keys does
not answer a wildcard request in one lookup.
